\section{Related Work}
\textit{Lexical URL detection.} Sahingoz~\textit{et al.}~\cite{sahingoz2019} released the EBBU-2017 corpus and showed that NLP-style URL features with tree ensembles reach high accuracy on that corpus. We reuse its raw URLs only; our 37 features are our own design and we do not attempt to reproduce their numbers.

\textit{Cross-dataset generalisation.} Rashid~\textit{et al.}~\cite{rashid2024} studied the generalisation of phishing URL classifiers across datasets and proposed unsupervised domain adaptation, which uses unlabelled target URLs during training; their public repository distributes EBBU, HP and ISCX files, from which we take HP. Our setting is strictly source-only; we do not replicate their algorithm or their splits. Kanon~\textit{et al.}~\cite{kanon2026} proposed cross-dataset feature stability selection: the intersection of top-$K$ Gini-important features across three datasets, two of which contribute a 15\% labelled sample of the external evaluation sets (pp.~9, 22 of~\cite{kanon2026}); splits are random at URL level. Misiek and Hyla~\cite{misiek2026} selected XGBoost features with criteria computed on importance across all datasets including the external ones and reported training/inference cost. Both studies therefore select features with information from the target population; neither purges shared domains. The closest idea to ours, stability of importance, is thus not new; what differs is \emph{where} stability is measured (held-out domain folds of the source only), the domain-disjoint protocol, the matched-budget contrast and the bidirectional design.

\textit{Selection bias and evaluation.} Feature selection performed outside the cross-validation loop biases error estimates~\cite{ambroise2002,varma2006}; we therefore nest selection inside source-only outer folds. Permutation importance~\cite{altmann2010} is computed on held-out domains and stability is borrowed in spirit from stability selection~\cite{meinshausen2010}. We report average precision rather than a trapezoidal PR-AUC~\cite{davis2006} and explain the actual tree models with TreeExplainer~\cite{lundberg2020}.
